%%%%%%%%%%%%%%%%%%%%%%%%%%%%%%%%%%%%%%%%%
% Medium Length Professional CV
% LaTeX Template
% Version 2.0 (8/5/13)
%
% This template has been downloaded from:
% http://www.LaTeXTemplates.com
%
% Original author:
% Trey Hunner (http://www.treyhunner.com/)
%
% Important note:
% This template requires the resume.cls file to be in the same directory as the
% .tex file. The resume.cls file provides the resume style used for structuring the
% document.
%
%%%%%%%%%%%%%%%%%%%%%%%%%%%%%%%%%%%%%%%%%

%----------------------------------------------------------------------------------------
%	PACKAGES AND OTHER DOCUMENT CONFIGURATIONS
%----------------------------------------------------------------------------------------

\documentclass{resume} % Use the custom resume.cls style
\usepackage[left=0.75in,top=0.6in,right=0.75in,bottom=0.6in]{geometry} % Document margins
\usepackage[utf8x]{inputenc}
\usepackage{eurosym}
\usepackage{enumitem} % Fine-grained control of list spacing
\usepackage{xcolor}
\usepackage{graphicx}
\usepackage[colorlinks=true]{hyperref}

\definecolor{linknavy}{RGB}{0,51,102}
\hypersetup{urlcolor=linknavy,linkcolor=linknavy,citecolor=linknavy}

\newcommand{\addphoto}[2]{%
  \smash{%
    \makebox[0pt][l]{%
      \raisebox{#1mm}{%
        \hspace{#2mm}\includegraphics[scale=0.40]{bz.jpg}%
      }%
     }%
  }%
}


\name{Alessio Brini, Ph.D} % Your name
\address{alessiobrini@gmail.com \quad alessio.brini@duke.edu}
\address{+1 (919) 904-5994}
\linkedin{\href{https://www.linkedin.com/in/alessiobrini/}{LinkedIn}}
\github{\href{https://github.com/Alessiobrini}{GitHub}}
\scholar{\href{https://scholar.google.com/citations?user=thrQtBcAAAAJ&hl=en}{Google Scholar}}
\ORCID{\href{https://orcid.org/0000-0003-1121-818X}{ORCID}}
\website{\href{https://alessiobrini.com}{alessiobrini.com}}

\begin{document}

%----------------------------------------------------------------------------------------
%	EDUCATION SECTION
%----------------------------------------------------------------------------------------
\begin{rSection}{Education}
{ Ph.D. at Scuola Normale Superiore } \hfill {11/2018 - 09/2022} \\
\emph{Mathematical Finance}. \hfill {\em Pisa, Italy}
\begin{list}{$\cdot$}{\leftmargin=0em}
\item Thesis: \textit{Reinforcement learning for sequential decision-making: a data-driven approach for finance}
\end{list}

{ Master's degree at Università degli studi di Firenze} \hfill {2016 - 2018} \\
\emph{Quantitative Finance and Risk Management} \hfill {\em Florence, Italy}
\begin{list}{$\cdot$}{\leftmargin=0em}
\item Thesis: ``Beyond the parameters' estimation burden: a Bayesian framework for option pricing.''
\end{list}

\vspace{0.5em}

{Bachelor's degree at Università degli studi di Firenze} \hfill {2012 - 2015} \\
\emph{Business and Economics}. \hfill {\em Florence, Italy}

\end{rSection}

%----------------------------------------------------------------------------------------
%	APPOINTMENTS SECTION
%----------------------------------------------------------------------------------------
\begin{rSection}{Academic Appointments}

{\bf Executive in Residence - Duke University - Pratt School of Engineering} \\
\smallskip
{07/01/2024 - Ongoing}, {\em Durham, NC, USA} \\
\begin{list}{$\cdot$}{\leftmargin=0em}
\item Leads the \textit{Digital Asset Research \& Engineering Collaborative (\href{https://darec.duke.edu/}{DAREC})} Lab.
\item My current role as a faculty member in FinTech continues my work in applying machine learning and data science methodologies to financial challenges, particularly in digital assets and decentralized finance (DeFi), while delivering advanced coursework and leading research projects in these areas.
\end{list}

{\bf Postdoctoral Researcher - Duke University - Pratt School of Engineering} \\
{01/23/2022 - 06/30/2024}, {\em Durham, NC, USA} \\
\begin{list}{$\cdot$}{\leftmargin=0em}
\item Applied machine learning and data science methodologies to address financial challenges, specializing in digital assets and decentralized finance (DeFi). Developed proficiency in diverse machine learning paradigms (supervised, unsupervised, reinforcement learning).
\item Delivered coursework on \emph{Machine Learning for Fintech}, \emph{Data Wrangling and Visualization with Python}, and \emph{Introduction to Statistics} for the Master of Engineering in Financial Technology program.
\item Directed FinTech Capstone Projects spanning Machine Learning, Data Science, and Digital Ledger Technologies.
\item Led summer research projects on DeFi and machine learning for graduate students.
\end{list}

{\bf Ph.D Researcher - Scuola Normale Superiore} \\
{11/2018 - 09/2022}, {\em Pisa, Italy} \\
\begin{list}{$\cdot$}{\leftmargin=0em}
\item Concentrated on the implementation of value-based and policy-based reinforcement learning algorithms to solve diverse economic and financial issues, including trading, portfolio selection, and best execution.
\end{list}

{\bf Actuarial Analyst - Milliman} \\ {03/2018 - 09/2018}, {\em Florence, Italy} \\
\begin{list}{$\cdot$}{\leftmargin=0em}
\item Utilized spreadsheets, databases, and other analytical tools to synthesize and evaluate data. Communicated findings to colleagues, project leads, and senior executives and compiled detailed written reports.
\end{list}
\end{rSection}

%----------------------------------------------------------------------------------------
%	PUBLICATIONS SECTION
%----------------------------------------------------------------------------------------
\begin{rSection}{Publications}

{\bf Peer-Reviewed Publications} \\[0.3em]
\begin{enumerate}[leftmargin=1.9em,itemsep=0.45em,topsep=0.15em,label=\arabic*.]
\item ``On Deep Reinforcement Learning for Dynamic Trading with PPO: Challenges and Future Directions'' (with Petter Kolm). \href{https://doi.org/10.3905/jfds.2026.001}{\emph{The Journal of Financial Data Science}}. \hfill 2026
\item ``Beyond Average Hive Performance: Tail Risk Measurement in Italian Apiculture with Honey-at-Risk'' (with Giacomo Toscano, Ginevra Virginia Lombardi, and Maria Elvira Mancino). \href{https://doi.org/10.1002/env.70106}{\emph{Environmetrics}}. \hfill 2026
\item ``Graph Neural Networks for Bridge Swap Link Prediction in Uniswap v3'' (with Qingran Zhou and Eric Liu). \href{https://doi.org/10.1145/3768292.3770392}{\emph{Proceedings of the 6th ACM International Conference on AI in Finance (ICAIF'25)}}. \hfill 2025
\item ``SpotV2Net: Multivariate Intraday Spot Volatility Forecasting via Vol-of-Vol-Informed Graph Attention Networks'' (with Giacomo Toscano). \href{https://doi.org/10.1016/j.ijforecast.2024.11.004}{\emph{International Journal of Forecasting}}. \hfill 2025
\item ``Modeling and Simulating a Startup Ecosystem Development with a Delayed Differential Equation'' (with Viviana Fanelli). \href{https://doi.org/10.1007/s00500-025-10565-4}{\emph{Soft Computing}}. \hfill 2025
\item ``Data-driven Derivative Hedging with Quadratic Variation Penalty'' (with Giacomo Domeniconi and Ali Fathi). \href{https://doi.org/10.1145/3677052.3698664}{\emph{Proceedings of the 5th ACM International Conference on AI in Finance (ICAIF'24)}}. \hfill 2024
\item ``Pricing Cryptocurrency Options with Machine Learning Regression for Handling Market Volatility'' (with Jimmie Lenz). \href{https://doi.org/10.1016/j.econmod.2024.106752}{\emph{Economic Modelling}}. \hfill 2024
\item ``A Comparison of Cryptocurrency Volatility-benchmarking New and Mature Asset Classes'' (with Jimmie Lenz). \href{https://doi.org/10.1186/s40854-024-00646-y}{\emph{Financial Innovation}}. \hfill 2024
\item ``Reinforcement Learning Policy Recommendation for Interbank Network Stability'' (with Gabriele Tedeschi and Daniele Tantari). \href{https://doi.org/10.1016/j.jfs.2023.101139}{\emph{Journal of Financial Stability}}. \hfill 2023
\item ``Deep Reinforcement Trading with Predictable Returns'' (with Daniele Tantari). \href{https://doi.org/10.1016/j.physa.2023.128901}{\emph{Physica A: Statistical Mechanics and its Applications}}. \hfill 2023
\item ``Assessing the resiliency of investors against cryptocurrency market crashes through the leverage effect'' (with Jimmie Lenz). \href{https://doi.org/10.1016/j.econlet.2022.110885}{\emph{Economics Letters}}. \hfill 2022
\end{enumerate}

\smallskip
{\bf Working Papers} \\[0.3em]
\begin{enumerate}[leftmargin=1.9em,itemsep=0.45em,topsep=0.15em,label=\arabic*.]
\item ``StableAML: Machine Learning for Behavioral Wallet Detection in Stablecoin Anti-Money Laundering on Ethereum'' (with Luciano Juvinski and Haochen Li). \href{https://arxiv.org/abs/2602.17842}{arXiv:2602.17842}. \hfill 2026
\item ``A Nonlinear Target-Factor Model with Attention Mechanism for Mixed-Frequency Data'' (with Ekaterina Seregina). \href{https://arxiv.org/abs/2601.16274}{arXiv:2601.16274}. \hfill 2026
\item ``Empirical Models of the Time Evolution of SPX Option Prices'' (with David Hsieh, Patrick Kuiper, Sean Moushegian, and David Ye). \href{https://arxiv.org/abs/2506.17511}{arXiv:2506.17511}. \hfill 2025
\item ``Improving DeFi Accessibility through Efficient Liquidity Provisioning with Deep Reinforcement Learning'' (with Haonan Xu). \href{https://arxiv.org/abs/2501.07508}{arXiv:2501.07508}. \hfill 2025
\item ``A Machine Learning Approach to Forecasting Honey Production with Tree-Based Methods'' (with Elisa Giovannini and Elia Smaniotto). \href{https://arxiv.org/abs/2304.01215}{arXiv:2304.01215}. \hfill 2023
\end{enumerate}

\smallskip
{\bf Other Articles and Lightly Refereed Publications} \\[0.3em]
\begin{itemize}[leftmargin=1.9em,itemsep=0.35em,topsep=0.15em,label=$\cdot$]
\item ``Crypto dynamics during market downturns: Another dotcom boom-bust cycle?'' (with Jimmie Lenz and Emma Rasiel). DAREC working paper. \hfill 2022
\item ``Bitcoin ETFs: Measuring the performance of this new market niche'' (with Jimmie Lenz). \href{https://papers.ssrn.com/sol3/papers.cfm?abstract_id=4157711}{SSRN 4157711}. \hfill 2022
\end{itemize}

\end{rSection}

%----------------------------------------------------------------------------------------
%	GRANT SECTION
%----------------------------------------------------------------------------------------
\begin{rSection}{Grants \& Funding}
\begin{list}{$\cdot$}{\leftmargin=0em}
\item Co-Investigator on the project ``Risk Management in Times of Unprecedented Geo-Political Volatility: A Machine-Learning Approach'', funded by the Institut Europlace de Finance (IEF), Institute Louis Bachelier (Grant amount: 10k \euro{}).
\item Collaborator in the ``BEEkeepers Weather Indexed INsurance'' project at the University of Florence, funded by the Italian Ministry of Agricultural, Food and Forestry Policies (Grant amount: 600k \euro{}).
\end{list}
\end{rSection}

%----------------------------------------------------------------------------------------
%	CONFERENCE SECTION
%----------------------------------------------------------------------------------------
\begin{rSection}{Conference Presentations \& Invited Talks}

\textbf{Quaint Quant Conference} - SMU Cox School of Business, Dallas, TX, April 10, 2026:
\begin{list}{$\cdot$}{\leftmargin=0em}
\item Contributed Talk: \emph{Attention-Based Mixed-Frequency Factor Models for Macroeconomic Forecasting}.
\end{list}

\textbf{Conference on Society-Centered AI} - Duke University, Durham, NC, February 12-14, 2026:
\begin{list}{$\cdot$}{\leftmargin=0em}
\item Invited Talk (New Faculty Research Spotlight): \emph{AI in Finance and Economics: From Data to Decisions that Matter}.
\end{list}

\textbf{ICAIF'25: 6th ACM International Conference on AI in Finance} - Singapore, November 16-17, 2025:
\begin{list}{$\cdot$}{\leftmargin=0em}
\item Contributed Talk: \emph{Graph Neural Networks for Bridge Swap Link Prediction in Uniswap v3}. Honorable mention for best paper award.
\end{list}

\textbf{AI for Social Impact: Bridging Innovations in Finance, Social Media, and Crime Prevention - AAAI Conference} - Philadelphia, March 3, 2025:
\begin{list}{$\cdot$}{\leftmargin=0em}
\item Contributed Talk: \emph{Improving DeFi Accessibility through Efficient Liquidity Provisioning with Deep Reinforcement Learning}. Best paper award (\href{https://sites.google.com/view/aaai-2025-ai-si/best-paper?authuser=0}{link}).
\end{list}

\textbf{ICAIF'24: 5th ACM International Conference on AI in Finance} - New York, November 14-17, 2024:
\begin{list}{$\cdot$}{\leftmargin=0em}
\item Poster Presentation: \emph{Data-driven Derivative Hedging with Quadratic Variation Penalty}.
\end{list}

\textbf{2024 NBER-NSF Time Series Conference} - Philadelphia, September 20-21, 2024:
\begin{list}{$\cdot$}{\leftmargin=0em}
\item Poster Presentation: \emph{SpotV2Net: Multivariate intraday spot volatility forecasting via vol-of-vol-informed graph attention networks}.
\end{list}

\textbf{Machine Learning of Dynamic Processes and Time Series Analysis} - Pisa, November 26-27, 2020:
\begin{list}{$\cdot$}{\leftmargin=0em}
\item Contributed Talk: \emph{Trading Mean Reversion with Reinforcement Learning}.
\end{list}

\textbf{13th European Financial Mathematics Summer School} - Vienna, August 31 - September 4, 2020:
\begin{list}{$\cdot$}{\leftmargin=0em}
\item Contributed Talk: \emph{Trading Mean Reversion with Reinforcement Learning}.
\end{list}

\textbf{Eastern European Machine Learning Summer School} - Krakow, July 1-9, 2020:
\begin{list}{$\cdot$}{\leftmargin=0em}
\item Poster Presentation: \emph{Trading Mean Reversion with Reinforcement Learning}.
\end{list}

\end{rSection}

%----------------------------------------------------------------------------------------
%	TEACHING SECTION
%----------------------------------------------------------------------------------------
\begin{rSection}{Teaching Experience}

\textbf{FINTECH 540: Machine Learning for FinTech - Pratt School of Engineering, Duke University}:

Explored the impact of machine learning on the rapidly evolving FinTech industry through both theoretical understanding and practical application, with extensive coding sessions utilizing real financial data.\\

\textbf{FINTECH/ECE 590: Data Wrangling and Visualization with Python - Pratt School of Engineering, Duke University}:

Leveraged Python for data querying, cleaning, manipulation, and visualization. Covered advanced topics such as Databases (MySQL, NoSQL), Web Scraping, Data Manipulation, and Data Collection.\\

\textbf{FINTECH 502: FinTech Capstones - Pratt School of Engineering, Duke University}:

Oversaw teams of students developing real-world technical solutions in the FinTech industry, with projects sponsored by industry professionals.\\

\textbf{FINTECH 520: Introduction to Statistics and Econometrics - Pratt School of Engineering, Duke University}:

Designed specifically for graduate students in Financial Technology, this course established fundamental statistical concepts and tools essential for advanced FinTech studies and research.\\

\textbf{Python for Data Science - University of Florence (Italy) - March-May, Sep-Dec, 2020}:

Provided an introduction to Python, covering the basics of the language structure and the main packages of the scientific stack for data analysis and extraction of meaningful information.\\

\end{rSection}

%----------------------------------------------------------------------------------------
%	ADVISING & MENTORING SECTION
%----------------------------------------------------------------------------------------
\begin{rSection}{Advising \& Mentoring}
\textbf{Graduate Students}:
\begin{list}{$\cdot$}{\leftmargin=0em}
\item Duccio Guerrini, University of Florence, 2020: ``Stock market indexes and machine learning algorithms. Replication of the S\&P500 index by the autoencoder neural network.''
\item Dmitry Guzairov, University of Florence, 2021: ``Hierarchical risk-parity portfolios: a clustering approach to the investment problem''
\item Niccolò Tarchi, University of Florence, 2021: ``Transformer-based approach for creating financial portfolios''
\item Mike Liu, Duke University, 2022: Independent Study on Asset Pricing and Financial Dashboard development
\item Alex Ilgenfritz, Duke University, 2023: ``Optimal Liquidity Provision in Uniswap v3 with Reinforcement Learning''
\item Haonan Xu, Duke University, 2024: ``Improving DeFi Accessibility through Efficient Liquidity Provisioning with Deep Reinforcement Learning''
\item Claire Zhou, Duke University, 2025: ``Graph Neural Networks for Bridge Swap Link Detection on Uniswap''
\item Longyi Hu, Duke University, 2025: ``Conformal Prediction for Finance and Economics.''
\item Luciano Silva, Duke University, 2025: ``A Scalable Framework for Stablecoin AML: Leveraging Supervised Learning Models for Behavioral Wallet Detection on Ethereum''
\end{list}
\end{rSection}

%----------------------------------------------------------------------------------------
%	SERVICE SECTION
%----------------------------------------------------------------------------------------
\begin{rSection}{Professional Service}

\textbf{Referee Service}:

Journal of Financial Econometrics, Annals of Operations Research, Expert Systems with Applications, Financial Innovation, Finance Research Letters, Applied Economics Letters, Scientific Reports, Applied Economics, Journal of Applied Economics, Discover Artificial Intelligence, Cogent Economics and Finance, AI Computer Science and Robotics Technology, International Conference in AI for Finance (ICAIF) 2024--2025, European Conference in AI for Finance (ECAIF) 2024--2025

\smallskip
\textbf{Conference Organization}:
\begin{list}{$\cdot$}{\leftmargin=0em}
\item Scientific Committee - Florence-Paris Workshop: Statistics of Random Processes and Its Applications to Financial Econometrics, July 2023, Florence, Italy
\item Organizer - Digital Assets at Duke (DA@D), January 2023 and 2024, Durham, NC, USA
\end{list}

\smallskip
\textbf{Departmental Service (Duke University, Pratt School of Engineering)}:
\begin{list}{$\cdot$}{\leftmargin=0em}
\item Assisting with recruitment and admissions review for the FinTech graduate program (2023-Ongoing).
\item Enhancing program outreach within the industry and research landscapes.
\end{list}

\end{rSection}

%----------------------------------------------------------------------------------------
%	TECHNICAL STRENGTHS SECTION
%----------------------------------------------------------------------------------------
\begin{rSection}{Technical Skills}
\begin{tabular}{ @{} >{\bfseries}l @{\hspace{6ex}} l }
Proficient Knowledge & Python, TensorFlow, PyTorch, SQL, Office suite (Excel) \\
Intermediate Knowledge & Docker, Matlab, R \\
Basic Knowledge & JavaScript, EViews \\
Languages & \textbf{Italian}: Native, \textbf{English}: Fully Proficient
\end{tabular}
\end{rSection}


\end{document}
